% Generated draft: requires mathematical and source review.
\documentclass[11pt]{article}
\usepackage[margin=25mm]{geometry}
\usepackage{fontspec}
\setmainfont{Latin Modern Roman}
\usepackage{amsmath,amssymb,mathtools,unicode-math}
\setmathfont{Latin Modern Math}
\usepackage{enumitem,needspace,xurl,hyperref}
\setlength{\parindent}{0pt}
\setlength{\parskip}{5pt}
\setlength{\emergencystretch}{3em}
\newcommand{\N}{\mathbb{N}}
\newcommand{\Z}{\mathbb{Z}}
\newcommand{\Q}{\mathbb{Q}}
\newcommand{\R}{\mathbb{R}}
\newcommand{\C}{\mathbb{C}}
\newcommand{\F}{\mathbb{F}}
\newcommand{\A}{\mathbb{A}}
\newcommand{\PP}{\mathbb{P}}
\newcommand{\E}{\mathbb{E}}
\newcommand{\eps}{\varepsilon}
\newcommand{\dd}{\,\mathrm{d}}
\begin{document}
\section*{UnsolvedMath ID: 11000097}
\textbf{Source problem code:} catalogue entry
\subsection{Definitions and mathematical statement}
Let \(\operatorname{Mod}_S\) be the mapping class group of the surface \(S\) in the source's convention. A pseudo-Anosov class preserves transverse measured foliations and scales them by reciprocal factors \(\lambda>1\) and \(\lambda^{-1}\). The intended existence question for a given \(S\) is
\[
 \exists N\triangleleft\operatorname{Mod}_S:\quad N\ne\{1\},
 \qquad\forall f\in N\setminus\{1\},\quad f\text{ is pseudo-Anosov}.
\]
Normality means \(hNh^{-1}=N\) for every \(h\in\operatorname{Mod}_S\). The nontriviality of \(N\) is necessary to avoid the vacuous trivial-subgroup answer, but is not explicitly stated in the source. No finite-generation, freeness or finite-index condition is requested.

The source credits K. Whittlesey [Whi] with examples for spheres with at least five punctures and for closed genus-two surfaces, noting the latter's connection with the six-punctured sphere. It emphasizes the still-intended higher-genus closed-surface cases and expresses pessimism about extending that particular proof method; the pessimism is a source opinion, not a mathematical impossibility result. The exact surface range, orientation/puncture conventions and whether the question seeks existence for every remaining type are not specified. The known source examples must not be used to turn the remaining question into the already answered assertion that some surface admits such a group.
\subsection{Short English statement}
For the intended surface, does its mapping class group contain a nontrivial normal subgroup all of whose nonidentity elements are pseudo-Anosov, beyond the examples reported for punctured spheres and genus two?
\subsection{Sources}
\begin{itemize}
\item N. Ivanov, chapter in \emph{Problems on Mapping Class Groups and Related Topics}, edited by Farb, Question I, supplied PDF page 82. \url{https://www.math.uchicago.edu/~farb/papers/mcgbook.pdf}. Recorded access 2026-07-29; extraction pdf-font-structure.
\item The source attributes independent posing to D. D. Long, J. D. McCarthy and R. C. Penner in the early 1980s, and says the author learned it from Penner in 1984. It cites K. Whittlesey [Whi] for the punctured-sphere and genus-two examples; the full bibliography entry is absent.
\item Local UnsolvedMath registry, \texttt{data/unsolvedmath/problems.json}, permanent ID \(11000097\), exact source code \texttt{catalogue entry}; supplied OPEN-TRIAGE summary is tentative about related results and supplies no exact later reference.
\end{itemize}
\end{document}
